\documentclass[9pt,t]{beamer}
% =====================================================================
% finance-beamer.tex — shared Beamer style for the LLM-in-Finance course
% =====================================================================
% Reusable slide style. Each deck \input's this immediately after
% \documentclass{beamer}. It provides:
%   * the Madrid theme + book colour palette (matches the printed book)
%   * two book-style framing blocks:
%       \begin{bigpicture} ... \end{bigpicture}   ("the bigger picture")
%       \begin{underhood}   ... \end{underhood}    ("under the hood")
%     These mirror the book's `context` and `deepdive` tcolorboxes so the
%     same intuition-vs-mechanism scaffold the reader meets in the book
%     reappears on the slides.
%   * a Python listings style for code frames
%   * appendix conventions: \appendixstart drops a divider and switches
%     to non-numbered backup frames so "extra / less central" material
%     lives after the main talk without inflating the slide count.
%
% Usage:
%   \documentclass{beamer}
%   % =====================================================================
% finance-beamer.tex — shared Beamer style for the LLM-in-Finance course
% =====================================================================
% Reusable slide style. Each deck \input's this immediately after
% \documentclass{beamer}. It provides:
%   * the Madrid theme + book colour palette (matches the printed book)
%   * two book-style framing blocks:
%       \begin{bigpicture} ... \end{bigpicture}   ("the bigger picture")
%       \begin{underhood}   ... \end{underhood}    ("under the hood")
%     These mirror the book's `context` and `deepdive` tcolorboxes so the
%     same intuition-vs-mechanism scaffold the reader meets in the book
%     reappears on the slides.
%   * a Python listings style for code frames
%   * appendix conventions: \appendixstart drops a divider and switches
%     to non-numbered backup frames so "extra / less central" material
%     lives after the main talk without inflating the slide count.
%
% Usage:
%   \documentclass{beamer}
%   % =====================================================================
% finance-beamer.tex — shared Beamer style for the LLM-in-Finance course
% =====================================================================
% Reusable slide style. Each deck \input's this immediately after
% \documentclass{beamer}. It provides:
%   * the Madrid theme + book colour palette (matches the printed book)
%   * two book-style framing blocks:
%       \begin{bigpicture} ... \end{bigpicture}   ("the bigger picture")
%       \begin{underhood}   ... \end{underhood}    ("under the hood")
%     These mirror the book's `context` and `deepdive` tcolorboxes so the
%     same intuition-vs-mechanism scaffold the reader meets in the book
%     reappears on the slides.
%   * a Python listings style for code frames
%   * appendix conventions: \appendixstart drops a divider and switches
%     to non-numbered backup frames so "extra / less central" material
%     lives after the main talk without inflating the slide count.
%
% Usage:
%   \documentclass{beamer}
%   \input{../_style/finance-beamer.tex}
%   \title{...}\subtitle{...}\author{...}\institute{...}\date{\today}
%   \begin{document} ... \end{document}
% =====================================================================

\usetheme{Madrid}
\usecolortheme{whale}
\setbeamertemplate{navigation symbols}{}
\setbeamertemplate{caption}[numbered]
\setbeamercovered{transparent}

% --- Density controls: keep dense, equation-heavy frames on one slide ---
% Slightly smaller body text + tighter lists/blocks. Frame titles and the
% Madrid chrome keep their normal size, so the deck still reads as Madrid,
% just with more vertical room for content.
\setbeamerfont{normal text}{size=\small}
\setbeamertemplate{itemize/enumerate body begin}{\small}
\setbeamertemplate{itemize/enumerate subbody begin}{\footnotesize}
% Tighter block spacing.
\setbeamertemplate{block begin}{%
  \par\vskip2pt%
  \begin{beamercolorbox}[colsep*=.75ex,leftskip=1ex,rightskip=1ex]{block title}%
    \usebeamerfont*{block title}\insertblocktitle%
  \end{beamercolorbox}%
  {\parskip0pt\vskip-1pt}%
  \begin{beamercolorbox}[colsep*=.75ex,vmode,leftskip=1ex,rightskip=1ex]{block body}%
  \vskip1pt}
\setbeamertemplate{block end}{\end{beamercolorbox}\vskip2pt}

\usepackage{amsmath, amssymb}
\usepackage{booktabs}
\usepackage{xcolor}
\usepackage{listings}
\usepackage[skins, breakable]{tcolorbox}

% --- Book colour palette (kept in sync with book/preamble.tex) ---
\definecolor{linkblue}{RGB}{14, 75, 140}
\definecolor{deepdivebg}{RGB}{235, 244, 255}
\definecolor{narrativebg}{RGB}{255, 252, 240}
\definecolor{narrativerule}{RGB}{180, 130, 40}
\definecolor{steelblue}{RGB}{70, 130, 180}
\definecolor{forestgreen}{RGB}{34, 139, 34}
\definecolor{crimson}{RGB}{220, 20, 60}
\definecolor{darkorange}{RGB}{255, 140, 0}
\definecolor{codebg}{RGB}{248, 248, 248}
\definecolor{coderule}{RGB}{210, 210, 210}
\definecolor{keywordblue}{RGB}{0, 100, 180}
\definecolor{commentgray}{RGB}{120, 120, 120}
\definecolor{stringgreen}{RGB}{30, 130, 30}

% Tie the Madrid structural colour to the book's link blue.
\setbeamercolor{structure}{fg=linkblue}
\setbeamercolor{title}{fg=white}
\setbeamercolor{frametitle}{fg=white}

% --- "the bigger picture" : narrative / intuition framing -----------
\newtcolorbox{bigpicture}{%
  enhanced, breakable, boxsep=2pt,
  colback  = narrativebg,
  colframe = narrativerule,
  boxrule  = 0pt, toprule = 1.4pt, bottomrule = 0.4pt,
  leftrule = 0pt, rightrule = 0pt, arc = 0pt,
  left = 6pt, right = 6pt, top = 4pt, bottom = 5pt,
  fonttitle = \footnotesize\scshape\color{narrativerule},
  title = {the bigger picture}, attach title to upper = {\par\smallskip},
}

% --- "under the hood" : the mechanism / technical detail ------------
\newtcolorbox{underhood}{%
  enhanced, breakable, boxsep=2pt,
  colback  = deepdivebg,
  colframe = linkblue,
  boxrule  = 0pt, toprule = 1.4pt, bottomrule = 0.4pt,
  leftrule = 0pt, rightrule = 0pt, arc = 0pt,
  left = 6pt, right = 6pt, top = 4pt, bottom = 5pt,
  fonttitle = \footnotesize\scshape\color{linkblue},
  title = {under the hood}, attach title to upper = {\par\smallskip},
}

% --- Python code style ---------------------------------------------
\lstset{
  basicstyle    = \ttfamily\footnotesize,
  keywordstyle  = \color{keywordblue}\bfseries,
  commentstyle  = \color{commentgray}\itshape,
  stringstyle   = \color{stringgreen},
  showstringspaces = false,
  language      = Python,
  breaklines    = true,
  backgroundcolor = \color{codebg},
  frame         = single,
  rulecolor     = \color{coderule},
  numbers       = none,
  aboveskip     = 4pt, belowskip = 4pt,
}

% --- Appendix conventions ------------------------------------------
% \appendixstart{Title} : begins the "extra material" appendix. Frames
% after it are not counted toward the main deck's frame total, keeping
% the core talk lean while preserving the deeper / optional content.
\newcommand{\appendixstart}[1]{%
  \appendix
  \setbeamertemplate{frametitle continuation}{}%
  \begin{frame}[noframenumbering]
    \begin{center}
      {\Large\scshape\color{linkblue} Appendix}\\[4pt]
      {\large #1}\\[10pt]
      {\footnotesize Extra / optional material — beyond the core lecture.}
    \end{center}
  \end{frame}
}
% Use \begin{frame}[noframenumbering]{...} for each appendix frame.

%   \title{...}\subtitle{...}\author{...}\institute{...}\date{\today}
%   \begin{document} ... \end{document}
% =====================================================================

\usetheme{Madrid}
\usecolortheme{whale}
\setbeamertemplate{navigation symbols}{}
\setbeamertemplate{caption}[numbered]
\setbeamercovered{transparent}

% --- Density controls: keep dense, equation-heavy frames on one slide ---
% Slightly smaller body text + tighter lists/blocks. Frame titles and the
% Madrid chrome keep their normal size, so the deck still reads as Madrid,
% just with more vertical room for content.
\setbeamerfont{normal text}{size=\small}
\setbeamertemplate{itemize/enumerate body begin}{\small}
\setbeamertemplate{itemize/enumerate subbody begin}{\footnotesize}
% Tighter block spacing.
\setbeamertemplate{block begin}{%
  \par\vskip2pt%
  \begin{beamercolorbox}[colsep*=.75ex,leftskip=1ex,rightskip=1ex]{block title}%
    \usebeamerfont*{block title}\insertblocktitle%
  \end{beamercolorbox}%
  {\parskip0pt\vskip-1pt}%
  \begin{beamercolorbox}[colsep*=.75ex,vmode,leftskip=1ex,rightskip=1ex]{block body}%
  \vskip1pt}
\setbeamertemplate{block end}{\end{beamercolorbox}\vskip2pt}

\usepackage{amsmath, amssymb}
\usepackage{booktabs}
\usepackage{xcolor}
\usepackage{listings}
\usepackage[skins, breakable]{tcolorbox}

% --- Book colour palette (kept in sync with book/preamble.tex) ---
\definecolor{linkblue}{RGB}{14, 75, 140}
\definecolor{deepdivebg}{RGB}{235, 244, 255}
\definecolor{narrativebg}{RGB}{255, 252, 240}
\definecolor{narrativerule}{RGB}{180, 130, 40}
\definecolor{steelblue}{RGB}{70, 130, 180}
\definecolor{forestgreen}{RGB}{34, 139, 34}
\definecolor{crimson}{RGB}{220, 20, 60}
\definecolor{darkorange}{RGB}{255, 140, 0}
\definecolor{codebg}{RGB}{248, 248, 248}
\definecolor{coderule}{RGB}{210, 210, 210}
\definecolor{keywordblue}{RGB}{0, 100, 180}
\definecolor{commentgray}{RGB}{120, 120, 120}
\definecolor{stringgreen}{RGB}{30, 130, 30}

% Tie the Madrid structural colour to the book's link blue.
\setbeamercolor{structure}{fg=linkblue}
\setbeamercolor{title}{fg=white}
\setbeamercolor{frametitle}{fg=white}

% --- "the bigger picture" : narrative / intuition framing -----------
\newtcolorbox{bigpicture}{%
  enhanced, breakable, boxsep=2pt,
  colback  = narrativebg,
  colframe = narrativerule,
  boxrule  = 0pt, toprule = 1.4pt, bottomrule = 0.4pt,
  leftrule = 0pt, rightrule = 0pt, arc = 0pt,
  left = 6pt, right = 6pt, top = 4pt, bottom = 5pt,
  fonttitle = \footnotesize\scshape\color{narrativerule},
  title = {the bigger picture}, attach title to upper = {\par\smallskip},
}

% --- "under the hood" : the mechanism / technical detail ------------
\newtcolorbox{underhood}{%
  enhanced, breakable, boxsep=2pt,
  colback  = deepdivebg,
  colframe = linkblue,
  boxrule  = 0pt, toprule = 1.4pt, bottomrule = 0.4pt,
  leftrule = 0pt, rightrule = 0pt, arc = 0pt,
  left = 6pt, right = 6pt, top = 4pt, bottom = 5pt,
  fonttitle = \footnotesize\scshape\color{linkblue},
  title = {under the hood}, attach title to upper = {\par\smallskip},
}

% --- Python code style ---------------------------------------------
\lstset{
  basicstyle    = \ttfamily\footnotesize,
  keywordstyle  = \color{keywordblue}\bfseries,
  commentstyle  = \color{commentgray}\itshape,
  stringstyle   = \color{stringgreen},
  showstringspaces = false,
  language      = Python,
  breaklines    = true,
  backgroundcolor = \color{codebg},
  frame         = single,
  rulecolor     = \color{coderule},
  numbers       = none,
  aboveskip     = 4pt, belowskip = 4pt,
}

% --- Appendix conventions ------------------------------------------
% \appendixstart{Title} : begins the "extra material" appendix. Frames
% after it are not counted toward the main deck's frame total, keeping
% the core talk lean while preserving the deeper / optional content.
\newcommand{\appendixstart}[1]{%
  \appendix
  \setbeamertemplate{frametitle continuation}{}%
  \begin{frame}[noframenumbering]
    \begin{center}
      {\Large\scshape\color{linkblue} Appendix}\\[4pt]
      {\large #1}\\[10pt]
      {\footnotesize Extra / optional material — beyond the core lecture.}
    \end{center}
  \end{frame}
}
% Use \begin{frame}[noframenumbering]{...} for each appendix frame.

%   \title{...}\subtitle{...}\author{...}\institute{...}\date{\today}
%   \begin{document} ... \end{document}
% =====================================================================

\usetheme{Madrid}
\usecolortheme{whale}
\setbeamertemplate{navigation symbols}{}
\setbeamertemplate{caption}[numbered]
\setbeamercovered{transparent}

% --- Density controls: keep dense, equation-heavy frames on one slide ---
% Slightly smaller body text + tighter lists/blocks. Frame titles and the
% Madrid chrome keep their normal size, so the deck still reads as Madrid,
% just with more vertical room for content.
\setbeamerfont{normal text}{size=\small}
\setbeamertemplate{itemize/enumerate body begin}{\small}
\setbeamertemplate{itemize/enumerate subbody begin}{\footnotesize}
% Tighter block spacing.
\setbeamertemplate{block begin}{%
  \par\vskip2pt%
  \begin{beamercolorbox}[colsep*=.75ex,leftskip=1ex,rightskip=1ex]{block title}%
    \usebeamerfont*{block title}\insertblocktitle%
  \end{beamercolorbox}%
  {\parskip0pt\vskip-1pt}%
  \begin{beamercolorbox}[colsep*=.75ex,vmode,leftskip=1ex,rightskip=1ex]{block body}%
  \vskip1pt}
\setbeamertemplate{block end}{\end{beamercolorbox}\vskip2pt}

\usepackage{amsmath, amssymb}
\usepackage{booktabs}
\usepackage{xcolor}
\usepackage{listings}
\usepackage[skins, breakable]{tcolorbox}

% --- Book colour palette (kept in sync with book/preamble.tex) ---
\definecolor{linkblue}{RGB}{14, 75, 140}
\definecolor{deepdivebg}{RGB}{235, 244, 255}
\definecolor{narrativebg}{RGB}{255, 252, 240}
\definecolor{narrativerule}{RGB}{180, 130, 40}
\definecolor{steelblue}{RGB}{70, 130, 180}
\definecolor{forestgreen}{RGB}{34, 139, 34}
\definecolor{crimson}{RGB}{220, 20, 60}
\definecolor{darkorange}{RGB}{255, 140, 0}
\definecolor{codebg}{RGB}{248, 248, 248}
\definecolor{coderule}{RGB}{210, 210, 210}
\definecolor{keywordblue}{RGB}{0, 100, 180}
\definecolor{commentgray}{RGB}{120, 120, 120}
\definecolor{stringgreen}{RGB}{30, 130, 30}

% Tie the Madrid structural colour to the book's link blue.
\setbeamercolor{structure}{fg=linkblue}
\setbeamercolor{title}{fg=white}
\setbeamercolor{frametitle}{fg=white}

% --- "the bigger picture" : narrative / intuition framing -----------
\newtcolorbox{bigpicture}{%
  enhanced, breakable, boxsep=2pt,
  colback  = narrativebg,
  colframe = narrativerule,
  boxrule  = 0pt, toprule = 1.4pt, bottomrule = 0.4pt,
  leftrule = 0pt, rightrule = 0pt, arc = 0pt,
  left = 6pt, right = 6pt, top = 4pt, bottom = 5pt,
  fonttitle = \footnotesize\scshape\color{narrativerule},
  title = {the bigger picture}, attach title to upper = {\par\smallskip},
}

% --- "under the hood" : the mechanism / technical detail ------------
\newtcolorbox{underhood}{%
  enhanced, breakable, boxsep=2pt,
  colback  = deepdivebg,
  colframe = linkblue,
  boxrule  = 0pt, toprule = 1.4pt, bottomrule = 0.4pt,
  leftrule = 0pt, rightrule = 0pt, arc = 0pt,
  left = 6pt, right = 6pt, top = 4pt, bottom = 5pt,
  fonttitle = \footnotesize\scshape\color{linkblue},
  title = {under the hood}, attach title to upper = {\par\smallskip},
}

% --- Python code style ---------------------------------------------
\lstset{
  basicstyle    = \ttfamily\footnotesize,
  keywordstyle  = \color{keywordblue}\bfseries,
  commentstyle  = \color{commentgray}\itshape,
  stringstyle   = \color{stringgreen},
  showstringspaces = false,
  language      = Python,
  breaklines    = true,
  backgroundcolor = \color{codebg},
  frame         = single,
  rulecolor     = \color{coderule},
  numbers       = none,
  aboveskip     = 4pt, belowskip = 4pt,
}

% --- Appendix conventions ------------------------------------------
% \appendixstart{Title} : begins the "extra material" appendix. Frames
% after it are not counted toward the main deck's frame total, keeping
% the core talk lean while preserving the deeper / optional content.
\newcommand{\appendixstart}[1]{%
  \appendix
  \setbeamertemplate{frametitle continuation}{}%
  \begin{frame}[noframenumbering]
    \begin{center}
      {\Large\scshape\color{linkblue} Appendix}\\[4pt]
      {\large #1}\\[10pt]
      {\footnotesize Extra / optional material — beyond the core lecture.}
    \end{center}
  \end{frame}
}
% Use \begin{frame}[noframenumbering]{...} for each appendix frame.


\title{Practical Session 08: Building \& Choosing Financial LLMs}
\subtitle{Hands-On: DAPT Budgeting, Hybrid Routing, and a FinBERT vs.\ LLM Bench}
\author{Juan F. Imbet}
\institute{Paris Dauphine -- PSL University}
\date{\today}

\begin{document}

\begin{frame}\titlepage\end{frame}

\begin{frame}{Outline}\tableofcontents\end{frame}

% ============================================================
\section{Session Overview \& Recap}
% ============================================================

\begin{frame}{Session Overview}
\textbf{Duration:} $\sim$1 hour (50 min work + 10 min debrief)

\medskip
\begin{block}{Timed agenda}
\begin{enumerate}
  \item \textbf{Recap} (8 min): the DAPT objective + the effective-cost formula.
  \item \textbf{Problem 1} (12 min): size a DAPT corpus and budget for a FinBERT-scale
        encoder.
  \item \textbf{Problem 2} (15 min): find the cost-optimal \textbf{hybrid routing}
        fraction $\rho$ for a filing classifier.
  \item \textbf{Case Study} (15 min): run a FinBERT-style vs.\ general-model
        sentiment comparison in Python (\texttt{code/notebooks/08-domain-specific-llms/}).
  \item \textbf{Discussion} (10 min): choose an architecture for three real briefs.
\end{enumerate}
\end{block}
\end{frame}

\begin{frame}{Recap: the Two Formulas You Will Use}
\begin{underhood}
\textbf{DAPT objective.} From general weights $\theta_0$, continue pre-training on a
domain corpus $\mathcal{D}_{\text{domain}}$:
\[
  \theta^\star = \arg\min_\theta \mathcal{L}_{\text{PT}}(\theta;\mathcal{D}_{\text{domain}}),
  \quad \mathcal{L}_{\text{PT}} \in \{\text{MLM},\ \text{CLM}\},
\]
reduced LR ($\sim 10\times$) to limit catastrophic forgetting.
\end{underhood}

\medskip
\footnotesize \textbf{Rule of thumb:} encoders need 1--5B domain tokens; the first
$\sim$1B yields most of the gain (diminishing returns).
\end{frame}

\begin{frame}{Recap: the Effective-Cost Formula}
\begin{underhood}
\textbf{Effective cost per correct output.}
\[
  C_{\text{eff}}(m,\rho) = \frac{c_m + \rho\,c_h}{p_m + \rho(1-p_m)},
\]
$c_m$ = model cost/query, $c_h$ = human-review cost, $p_m$ = model correctness,
$\rho$ = fraction routed to review. Constraint: $p_m + \rho(1-p_m) \ge p^\star$.
\end{underhood}

\medskip
\begin{block}{How to read it}
The numerator is the cost you pay per query (model always, human only on the routed
fraction $\rho$); the denominator is the accuracy you achieve. Routing more to humans
raises both --- the optimisation finds the smallest $\rho$ that still meets $p^\star$.
\end{block}
\end{frame}

% ============================================================
\section{Guided Problems}
% ============================================================

\begin{frame}{Problem 1: Sizing a DAPT Corpus}
\textbf{Scenario.} You will DAPT a BERT-base encoder (110M params) for SEC-filing NER,
then fine-tune. You can draw from:
\begin{itemize}
  \item SEC EDGAR ($\sim$50B tokens, open)
  \item Reuters TRC2 ($\sim$2B, licensed)
  \item General replay data (Pile, $\sim$unlimited)
\end{itemize}

\medskip
\textbf{Tasks.}
\begin{enumerate}
  \item[\textbf{A.}] Given diminishing returns, what domain-token budget do you target,
        and why not just use all 50B?
  \item[\textbf{B.}] You worry about catastrophic forgetting. Propose a \emph{replay}
        ratio (general:domain) and justify it.
  \item[\textbf{C.}] Should you set the DAPT learning rate above or below the original
        pre-training LR? By roughly how much, and what failure does this prevent?
  \item[\textbf{D.}] Would TAPT help here? What additional data would it require?
\end{enumerate}
\end{frame}

\begin{frame}{Problem 2: Optimal Hybrid Routing}
\textbf{Scenario.} A FinBERT classifier costs $c_m=\$0.0002$/query at $p_m=0.88$.
Hard queries route to a human at $c_h=\$2.00$/query (assume review fixes all errors).
Target accuracy $p^\star = 0.95$.

\medskip
\textbf{Tasks.}
\begin{enumerate}
  \item[\textbf{A.}] Write the achieved accuracy as a function of $\rho$ and find the
        \emph{minimum} $\rho$ that meets $p^\star$.
  \item[\textbf{B.}] Evaluate $C_{\text{eff}}$ at that $\rho$.
  \item[\textbf{C.}] Compare against routing \emph{everything} to the human ($\rho=1$).
        What is the cost ratio?
  \item[\textbf{D.}] Qualitatively: if $p_m$ rose to $0.95$ via DAPT, what happens to
        the required $\rho$ and to $C_{\text{eff}}$?
\end{enumerate}

\medskip
\begin{block}{Hint}
Achieved accuracy $= p_m + \rho(1-p_m)$ because review recovers the $(1-p_m)$ errors
on the routed fraction.
\end{block}
\end{frame}

% ============================================================
\section{Python Case Study}
% ============================================================

\begin{frame}{Case Study: FinBERT-style vs.\ General Sentiment}
\textbf{Goal.} Demonstrate \emph{polarity inversion} --- the single clearest reason
finance needs its own models --- on a few labelled financial sentences.

\medskip
\textbf{Notebook:} \texttt{code/notebooks/08-domain-specific-llms/demo.ipynb}

\medskip
\textbf{Stack:}
\begin{itemize}
  \item \texttt{transformers} --- load \texttt{ProsusAI/finbert} (the Araci FinBERT)
        and a general sentiment model.
  \item A handful of finance sentences where general English polarity is misleading.
\end{itemize}

\medskip
\begin{alertblock}{Prerequisites}
\texttt{pip install transformers torch}. First run downloads model weights
($\sim$440MB for FinBERT). No GPU required for a few sentences.
\end{alertblock}
\end{frame}

\begin{frame}[fragile]{Case Study: The Code (1/2) --- Setup}
\begin{lstlisting}
from transformers import pipeline

# Domain-specific encoder (Araci FinBERT) vs. general sentiment
finbert = pipeline("text-classification", model="ProsusAI/finbert")
general = pipeline("sentiment-analysis")  # default SST-2 model

SENTENCES = [
    "The firm took on additional liability to fund the acquisition.",
    "We are cautiously optimistic but flag meaningful headwinds.",
    "Leverage increased to 3.2x following the debt issuance.",
    "Unprecedented demand drove revenue above guidance.",
]
\end{lstlisting}

\medskip
\footnotesize Two pipelines, four finance sentences chosen so general-English polarity
is misleading. The comparison loop is on the next slide.
\end{frame}

\begin{frame}[fragile]{Case Study: The Code (2/2) --- Comparison Loop}
\begin{lstlisting}
for s in SENTENCES:
    f = finbert(s)[0]
    g = general(s)[0]
    print(f"{s[:48]:48s} | FinBERT={f['label']:8s} "
          f"| General={g['label']}")
\end{lstlisting}

\medskip
\footnotesize Run it, then fill in the table on the next slide.
\end{frame}

\begin{frame}{Case Study: Diagnostic Questions}
\textbf{Q1.} For ``additional \emph{liability}'' and ``\emph{leverage} increased'',
what does the general model predict vs.\ FinBERT? Which is right \emph{for a financial
reader}, and why does the general model err?

\medskip
\textbf{Q2.} ``\emph{Cautiously optimistic} but \dots headwinds'' --- which label is
appropriate? Connect this to the FinBERT-tone earnings-call example from the lecture.

\medskip
\textbf{Q3.} ``\emph{Unprecedented} demand \dots'' is genuinely positive here. Does
FinBERT over-correct? What does this tell you about \emph{context}-dependence of
financial polarity?

\medskip
\begin{block}{Q4 --- tie to the chapter}
Of the four general-model failure modes (polarity inversion, numerical hallucination,
table-parsing, regulatory ambiguity), which one does this experiment isolate? Which
\emph{cannot} be diagnosed by a sentence-level sentiment test, and why?
\end{block}
\end{frame}

% ============================================================
\section{Discussion}
% ============================================================

\begin{frame}{Discussion: Pick an Architecture for Each Brief}
For each, state the model family + training strategy and justify on accuracy, cost,
latency, governance, and traceability.

\medskip
\begin{enumerate}
  \item A bank must label \textbf{50{,}000 headlines/day} as material / non-material
        for insider-trading surveillance. Latency $<$50ms; 10k labelled examples.
  \item A research desk needs \textbf{hawkish/dovish} scoring of every new FOMC
        statement, with the precise monetary-policy vocabulary.
  \item A compliance team must answer questions about \textbf{any 10-K} in a
        2{,}000-name portfolio, citing the exact source paragraph for audit.
\end{enumerate}

\medskip
\begin{block}{Frame your answer}
Encoder vs.\ decoder? From-scratch vs.\ DAPT vs.\ fine-tune-only? Single model vs.\
hybrid routing vs.\ RAG? Where does the LM dictionary / register argument bite?
\end{block}
\end{frame}

% ============================================================
% SOLUTIONS
% ============================================================

\begin{frame}{Solution: Problem 1 --- DAPT Corpus}
\begin{block}{A --- Budget}
Target $\sim$\textbf{3--5B} EDGAR tokens. Most of the gain is in the first $\sim$1B;
beyond $\sim$5B the marginal downstream improvement for a 110M encoder is small and
not worth the compute or the added risk of forgetting.
\end{block}
\textbf{B --- Replay.} Mix in $\sim$5--10\% general (Pile) tokens. Enough to preserve
broad language ability, small enough that the domain still dominates the update.

\medskip
\textbf{C --- Learning rate.} \emph{Below} the original PT rate, by $\sim 10\times$.
This prevents catastrophic forgetting --- a large LR would overwrite the general
representations DAPT is meant to \emph{reuse}.

\medskip
\textbf{D --- TAPT.} Yes: a further continued-PT pass on the \emph{unlabelled} text of
the NER dataset (same distribution as the labelled set), before supervised fine-tuning.
Requires task-distribution unlabelled data, which here we have.
\end{frame}

\begin{frame}{Solution: Problem 2 --- Optimal Routing}
\textbf{A --- Required $\rho$.} Achieved accuracy $= p_m + \rho(1-p_m) = 0.88 + 0.12\rho$.
Set $\ge 0.95$:
\[
  0.88 + 0.12\rho \ge 0.95 \;\Rightarrow\; \rho \ge \tfrac{0.07}{0.12} \approx \mathbf{0.583}.
\]

\textbf{B --- $C_{\text{eff}}$ at $\rho=0.583$:}
\[
  C_{\text{eff}} = \frac{0.0002 + 0.583\cdot 2.00}{0.88 + 0.583\cdot 0.12}
  = \frac{1.1662}{0.95} \approx \mathbf{\$1.23}\ \text{per correct output.}
\]

\textbf{C --- All-human ($\rho=1$):} $C_{\text{eff}} = (0.0002+2.00)/1.0 \approx \$2.00$.
Hybrid is $\approx$\textbf{1.6$\times$ cheaper} at the \emph{same} 95\% target.

\medskip
\textbf{D --- If DAPT lifts $p_m$ to 0.95.} Required $\rho = 0$ (target already met),
so $C_{\text{eff}} \approx c_m/p_m = 0.0002/0.95 \approx \$0.0002$ --- four orders of
magnitude cheaper. \emph{Investing in the model collapses the human-review bill.}
\end{frame}

\begin{frame}{Wrap-Up: What This Session Showed}
\begin{itemize}
  \item \textbf{Problem 1:} DAPT is a budgeting exercise --- a few billion domain
        tokens, a reduced LR, and a replay ratio buy most of the gain. More data is
        not always better.
  \item \textbf{Problem 2:} hybrid routing meets an accuracy target at a fraction of
        all-human cost --- and the value of \emph{model} improvement (DAPT raising
        $p_m$) is dominated by the human-review savings it unlocks.
  \item \textbf{Case study:} polarity inversion is real and visible in four sentences
        --- the cleanest empirical motivation for domain-specific models.
  \item \textbf{Discussion:} architecture choice follows the task's label space,
        vocabulary sensitivity, latency, and \emph{auditability} requirements.
\end{itemize}

\medskip
\begin{block}{Next}
Financial NLP \& sentiment in depth: fine-tuning FinBERT end-to-end and benchmarking
it against zero-shot frontier models.
\end{block}
\end{frame}

% ============================================================
% APPENDIX
% ============================================================
\appendixstart{Stretch Problems}

\begin{frame}[noframenumbering]{Appendix --- Stretch Problem 1: Contamination}
A vendor claims their 70B model scores 0.97 F1 on Financial PhraseBank --- above any
published FinBERT.

\medskip
\textbf{Tasks.}
\begin{enumerate}
  \item Why is this number \emph{not} immediately trustworthy?
  \item Design a contamination-controlled evaluation: what temporal property must the
        test set satisfy relative to the model's knowledge cutoff?
  \item PhraseBank predates almost every frontier model. What kind of held-out set
        would you construct instead, and from what source?
\end{enumerate}

\medskip
\footnotesize \textit{Key idea:} a benchmark in the pre-training corpus measures
memorisation, not generalisation. Require test data postdating the cutoff.
\end{frame}

\begin{frame}[noframenumbering]{Appendix --- Stretch Problem 2: RAG vs.\ Bigger Model}
A 7B decoder + EDGAR RAG vs.\ a 50B ungrounded decoder, both for ``extract the FY
revenue from this 10-K''.

\medskip
\textbf{Tasks.}
\begin{enumerate}
  \item On which axis does RAG win even if the 50B model is slightly more accurate
        ungrounded? (Hint: SR 11-7.)
  \item Sketch why grounding reduces \emph{numerical} hallucination specifically.
  \item Give one task where the 50B model would still be preferable, and say why
        retrieval does not help there.
\end{enumerate}

\medskip
\footnotesize \textit{Expected:} RAG wins on traceability + cost; grounding pins
figures to retrieved text; the big model still wins on open-ended reasoning that no
single passage supports.
\end{frame}

\end{document}
